\chapter*{Preface}
\addcontentsline{toc}{chapter}{Preface}
\markboth{Preface}{Preface}

This handbook is a complete engineering account of one real system: the
recipe from Mia's AI Lab that serves \textbf{DeepSeek-V4-Flash} --- a frontier-class,
million-token-context, mixture-of-experts language model with native vision
input --- on \textbf{two NVIDIA DGX Spark} desktop machines connected by a
single 200\,Gb RoCE cable. The system is public: it lives in the repository
\repofile{MiaAI-Lab/DeepSeek-v4-Flash-DSpark-2x-DGX-Spark}, together with
every patch, benchmark, retraction, and dated changelog entry this book draws
on. Nothing here is invented; where the repository itself marks evidence as
partial or pending, this book says so.

Why write it down at all? Because the hard part of this work is no longer
getting the hardware. Two DGX Sparks fit on a desk, and the weights are a
download away. The hard part is knowing which of a hundred flags actually
matter, which innocuous-looking defaults encode an incident somebody already
survived, and what a given measured number is evidence \emph{of}. That
knowledge exists here in full --- but it exists as 42 patch sources, 66
scripts, a changelog, and several hundred issue comments. That is the shape
evidence takes, not the shape understanding takes. Reading it in order of
commit is not the same as learning it.

This book is the assembly. Same facts, reordered until they teach: mechanism
before configuration, the incident before the flag that now prevents it,
and every number still carrying the date and lane it was measured on. If it
succeeds, a reader finishes able to reason about a serving stack they have
never seen --- because the constraints that shaped this one are not
peculiar to it.

The book is organized in two parts. \textbf{Part I} (Chapters 1--13) is the
DeepSeek recipe, told end to end. \textbf{Part II} (Chapters 14--17) studies
the lab's sibling recipe, \repofile{MiaAI-Lab/GLM-5.3-Flash-EXL3-2x-DGX-Sparks},
which serves \textbf{GLM-5.3-Flash} on the same two-Spark kit but makes the
opposite bet at nearly every layer: community-quantized EXL3 weights instead
of native NVFP4, a DFlash2 drafter instead of DSpark, and purpose-built CUDA
kernels where Part I leaned on the stock stack. Part II covers only what
Part I does not --- the concepts that appear when the same hardware is pushed
down a different road --- so the two parts together form one comparative
study rather than two summaries.

\section*{Who this book is for}

The primary audience is \textbf{AI architect engineers} and \textbf{inference
engineers}: people who design, deploy, and debug LLM serving systems and want
to see, at full resolution, what it takes to run a frontier model outside a
hyperscaler datacenter. The secondary audience is any strong technical reader.
No prior knowledge of vLLM internals, RoCE networking, or speculative decoding
is assumed --- each is taught from its mechanism before its configuration.

Why study this particular system? Three reasons. First, it is
\emph{small enough to understand completely}: two nodes, one model, one
repository --- yet it exercises nearly every hard problem in modern inference:
tensor parallelism over a real fabric, paged and quantized KV caches, sparse
attention, speculative decoding under continuous batching, CUDA-graph capture,
structured output grammars, multimodal input, and multi-week production
operation. Second, it is \emph{unusually honest}: benchmark tables carry
dates and noise bands, defaults are flipped back when new evidence arrives,
and at least one published safety claim is formally retracted in-tree. Third,
its engineering method --- digest-pinned images modified by fail-closed,
source-locked boot-time patches --- is a transferable discipline that this
book treats as a first-class subject.

\section*{Prerequisites}

The book assumes a practitioner's working knowledge of transformer inference
rather than a researcher's: what a token is, why a KV cache exists, and why
prefill and decode have different cost profiles. It assumes comfort on a
Linux host --- shell, containers, service supervision --- and enough Python
to read a patch, since patch sources are quoted throughout. It does not
assume any ability to derive attention, train a model, or write CUDA.

Everything specific to this system is taught before it is used. vLLM's
scheduler and paged KV allocator, multi-head latent attention, sparse
attention indexing, speculative decoding, CUDA-graph capture, RoCE and NCCL
device selection, and Part~II's quantization formats each get their mechanism
explained before any flag that configures them appears. A reader who has
never opened a serving engine's source can follow the argument end to end.

Reproducing the system, as opposed to reading about it, needs the hardware it
was measured on:

\begin{itemize}
  \item \textbf{Two NVIDIA DGX Spark machines} (GB10, 128\,GB unified
    LPDDR5X each). One node cannot hold the checkpoint: roughly 79\,GiB of
    weights land on each rank, leaving 15--17\,GiB for KV.
  \item \textbf{A direct ConnectX-7 link} --- the 200\,Gb QSFP ports cabled
    between the two machines, with RoCEv2 reachable on both, plus a
    management LAN for SSH and client traffic.
  \item \textbf{Docker with Compose on both hosts}, passwordless SSH from
    head to worker, and a Hugging Face token for the gated checkpoint.
\end{itemize}

That kit is where these numbers were measured, not what the book requires to
be useful. The cost models, the failure modes, and the method transfer to any
stack facing the same constraints --- a bandwidth-bound accelerator, a fabric
between ranks, a pinned image that has to be patched rather than forked ---
and the chapters are written so that each argument stands on its own
reasoning, with the hardware supplying the evidence rather than the premise.

\section*{Using this book to teach}

The material is arranged for a course, not only for reference. Every chapter
opens on a concrete failure or measurement, teaches the mechanism that
explains it, and closes with takeaways that state a transferable rule. The
callout boxes keep the four registers separate --- context
(\textbf{Note}), footguns (\textbf{Hazard}), principles (\textbf{Field
lesson}), narrative (\textbf{War story}) --- so any of them can be lifted for
a slide without dragging the surrounding prose along.

Seven modules group the chapters into teachable sessions:

\begin{center}\small
  \begin{tabular}{@{}ll>{\raggedright\arraybackslash}p{68mm}@{}}
    \toprule
    Module & Chapters & The question it answers \\
    \midrule
    1. Foundations & 1--3 & What is the machine, and what is the model? \\
    2. Configuration and fabric & 4--5 & How is it launched, and how do the ranks talk? \\
    3. Change control & 6--7 & How do you modify a pinned system without breaking it? \\
    4. Throughput & 8--9 & Where does decode time go, and what actually moves it? \\
    5. Evidence and operations & 10--11 & How do you know a number is real, and how is this run daily? \\
    6. Scale and synthesis & 12--13 & What lies beyond two nodes, and what generalizes? \\
    7. Comparative study & 14--17 & What changes when the same hardware takes the opposite bet? \\
    \bottomrule
  \end{tabular}
\end{center}

Two properties make this corpus unusually teachable. First, the evidence is
public: every claim traces back to a patch, a script, or a dated benchmark
file a student can open in a browser. Second, the sources record their own
mistakes --- a retracted safety matrix, a default flipped and flipped back, a
benchmark disowned as an artifact of a pathological prompt. Systems that
publish only their successes teach facts; systems that publish their
reversals teach judgment. Exercises follow almost for free: hand students the
symptom and the repository, and ask them to reach the diagnosis the chapter
reaches --- then compare their reasoning to the one the maintainers actually
used, which is also on record.

\section*{A living document}

This is an edition of a system still in motion, and the book is versioned
accordingly. Two revisions are planned rather than merely hoped for.

The larger one is \textbf{four Sparks}. \Cref{ch:transplants} closes on
\repofile{start-tp4.sh}, a TP=4 launcher shipped honestly unmeasured --- its
own header reads ``untested here (no 4-Spark kit).'' That is the frontier of
this material, and the one place where the book documents a configuration
nobody has yet run. When a four-node kit exists, the TP=4 lane will get what
\Cref{ch:scaling} gave the third Spark: a measured ledger, a fabric census,
and a verdict. The interesting question is already framed. Adding a third
node cut the bytes each rank must stream and moved the end metric by nothing,
because fixed per-step cost ate the entire bandwidth gain. Whether a fourth
node repeats that result, or finds a different wall, is not something the
byte model of \Cref{ch:perf} can settle by arithmetic. It has to be measured.

The second is \textbf{heavier models}. Both recipes here serve at the edge of
what two 128\,GB nodes hold. A larger checkpoint, or the same one at higher
precision, changes the weights-versus-KV split that \Cref{ch:profile} treats
as a fixed premise --- and every capacity figure in this book sits downstream
of that split. Those chapters will be re-derived rather than patched.

Until then, the honest scope of this edition is two nodes. Where the book
looks past that boundary it says so plainly, because
\Cref{ch:transplants} is explicit that a claim about hardware nobody has
measured is not a claim its own authors are willing to make. This book holds
itself to the same standard.

\section*{How this book was built}

The text was produced from an exhaustive read of the repository at commit
\code{f5665e8} (September 2026): all documentation, the complete changelog,
all 42 patch sources (14{,}443 lines), all 66 scripts (19{,}154 lines), the
launcher, the Compose files, the test suites, the source overlays, and every
results file. Measured numbers are quoted from the repository's own dated
records and are labeled with their measurement date and lane, because the
repository itself insists on that discipline. Part II was produced the same
way from an exhaustive read of the GLM repository (September 2026): both
launchers, all 21 overlay sources including the CUDA kernels, the full test
and bench suites, the abliteration pipeline, the design documents, and all
125 commits of its history.

\section*{Conventions}

File paths appear as \repofile{docker-compose.dspark.yml}, environment
variables as \envvar{DSPARK_MAX_INFLIGHT_PREFILLS}, code identifiers as
\code{max_num_partial_prefills}, and GitHub issues in the repository's tracker
as \ghissue{27}. Throughput is written in tokens per second (\tokspersec{}).
Four kinds of callout boxes recur:

\begin{hbnote}
\textbf{Note} boxes add context or point to related sections.
\end{hbnote}

\begin{hbwarning}
\textbf{Hazard} boxes mark footguns that have actually fired --- configuration
traps, silent failure modes, and irreversible operations.
\end{hbwarning}

\begin{hblesson}
\textbf{Field lesson} boxes distill a transferable engineering principle from
a concrete incident in this system.
\end{hblesson}

\begin{hbwarstory}
\textbf{War story} boxes recount a real incident, in the order it unfolded,
including the wrong turns.
\end{hbwarstory}

\section*{Acknowledgements}

This book exists because \textbf{Mia's AI Lab} does its engineering in the
open. Both recipes, and every number quoted here, are public because the lab
publishes its patches, its benchmarks, its dated changelogs --- and, more
unusually, its retractions and its unexplained results. A handbook of this
kind cannot be written about a system whose evidence is private. My thanks to
the lab for the work and for the openness:
\href{https://github.com/MiaAI-Lab}{github.com/MiaAI-Lab} and
\href{https://x.com/MiaAI_lab}{@MiaAI\_lab}.

The system documented here is the work of a community: Mia's AI Lab, who
packaged and operate the recipe; \emph{drowzeys} (``Keys''), whose concurrency
patches made batched DSpark correct; the Anemll project, whose prebuilt GB10
vLLM image is the runtime foundation; Rafael Caricio and tonyd2wild, whose
replicas established the reproduction baseline; \emph{localaiguyy}, whose
attention-group padding enabled the three-node lane; and the many issue
reporters whose numbered bugs structure half of this book. Their full credits
live in \repofile{CREDITS.md}.
